\documentclass[11pt]{article}
\usepackage[a4paper,margin=1in]{geometry}
\usepackage{amsmath,amssymb,amsthm,microtype}
\newtheorem{theorem}{Theorem}[section]
\newtheorem{corollary}[theorem]{Corollary}
\theoremstyle{remark}\newtheorem{remark}[theorem]{Remark}
\title{Root-of-Unity Packets that Isolate a Single Power Moment}
\author{Bin Wang}\date{July 2026}
\begin{document}\maketitle
\begin{abstract}
We construct a genuine finite normal spectrum that changes one prescribed
power moment without changing any lower moment.  The packet consists of a
complete root set of a real monomial equation, repeated with integer
multiplicity.  Its entire moment sequence, radius, rank, and squared mass are
explicit.  This supplies the triangular atom needed for exact finite-prefix
spectral realization, while making no identification with the actual noisy
operator.
\end{abstract}

\section{Packet definition}
Fix $d,L\ge1$ and $w\in\mathbb R\setminus\{0\}$.  Let
\[
 \mathcal P_d(w,L)=
 \{\mu:\mu^d=w/(dL)\},
\]
with every root repeated $L$ times.  Since the defining polynomial has real
coefficients, the multiset is closed under complex conjugation.

\begin{theorem}[moment isolation]
The packet has rank $dL$, radius
\[
 r_d=(|w|/(dL))^{1/d},
\]
and power sums
\[
 p_n(\mathcal P_d)=0\quad(d\nmid n),
 \qquad
 p_{md}(\mathcal P_d)=\frac{w^m}{(dL)^{m-1}}.
\]
In particular, $p_n=0$ for $n<d$ and $p_d=w$.  Its squared spectral mass is
\[
 M_2=dLr_d^2.
\]
\end{theorem}
\begin{proof}
Write the roots as $r_de^{i(\theta+2\pi j)/d}$.  The roots-of-unity sum
vanishes unless $d$ divides $n$.  At $n=md$, summing and multiplying by $L$
gives $Ld(w/(dL))^m$, which is the stated formula.  The geometric quantities
are immediate.
\end{proof}

\begin{corollary}[modulus cap]
For $q>0$, if
\[
 L\ge\frac{|w|}{dq^d},
\]
then every packet root satisfies $|\mu|\le q$.
\end{corollary}

\begin{remark}
The packet is an unweighted finite spectrum with integer multiplicities, not
merely a positive moment measure.  It is nevertheless synthetic: no theorem
identifies it with eigenvalues of the noisy transfer operator.  Gates A--E
remain false/open.
\end{remark}
\end{document}
